\section{实验任务与要求}
实验任务包括：在前车尾部粘贴指定 ID 的 ALVAR AR 标记；启动 Turtlebot2 底盘、Kinect 相机和 AR 识别节点；订阅 \texttt{/ar\_pose\_marker} 获取目标标记相对位姿；根据前后距离误差和左右偏差生成速度指令；通过速度门控限制危险速度并实现前方障碍停车；最后通过修改 \texttt{target\_id} 或 \texttt{marker\_id} 使机器人跟随不同编号的 AR 标记。

实验要求后车不能盲目运动；目标丢失、急停信号有效、扫描数据异常或前方距离过近时应停车。调试时应先单车低速验证，再扩展到多车队列。不同车辆应使用不同 ROS 命名空间和不同 TF 坐标系，避免话题和坐标冲突。
